% Multi-Source Knowledge Graph RAG for Medical Diagnosis
% v17 - 15-page version
% Target: Information Fusion (Elsevier)

\documentclass[preprint,12pt]{elsarticle}

\usepackage{amsmath,amssymb,amsfonts}
\usepackage{graphicx}
\usepackage{booktabs}
\usepackage{multirow}
\usepackage{hyperref}
\usepackage{algorithm}
\usepackage{algorithmic}
\usepackage{xcolor}
\usepackage{url}

\journal{Information Fusion}

\begin{document}

\begin{frontmatter}

\title{Multi-Source Knowledge Graph Retrieval-Augmented Generation\\as a Second Brain for LLM-Based Medical Diagnosis}

\author[addr1]{First Author}
\author[addr1]{Second Author}
\author[addr1]{Third Author}

\address[addr1]{Department of Computer Science, University, Country}

\begin{abstract}
Large Language Models (LLMs) suffer from hallucination and knowledge gaps in medical domains. We propose a multi-source Retrieval-Augmented Generation (RAG) framework functioning as a ``Second Brain'' for LLM-based differential diagnosis, integrating three complementary retrieval layers: (1)~dense vector retrieval over medical knowledge bases, (2)~sparse lexical retrieval enhanced with Positive Pointwise Mutual Information (PPMI) co-occurrence statistics, and (3)~a multi-source knowledge graph combining ontological structure with data-driven co-occurrence patterns. On DDXPlus (49~diseases, $n{=}200$), our pipeline achieves 95.5\% Top-1 and 98.5\% Top-5 accuracy with zero hallucinations, versus the unaugmented LLM's 32.5\% Top-1 and 59.8\% hallucination rate. On Symptom2Disease (24~diseases, $n{=}320$) with a domain-specific knowledge graph, it achieves 90.3\% Top-1 and 98.8\% Top-5 with zero hallucinations versus 75.9\% for the baseline. A GPT-4o ablation shows both models converge to identical performance under RAG, demonstrating that retrieval compensates for base model variation. Each retrieval layer contributes progressively, and our results demonstrate that structured multi-source RAG effectively eliminates hallucination while achieving near-perfect diagnostic accuracy.
\end{abstract}

\begin{keyword}
Retrieval-Augmented Generation \sep Knowledge Graph \sep Medical Diagnosis \sep Large Language Models \sep Hallucination Mitigation \sep Second Brain
\end{keyword}

\end{frontmatter}

%% ============================================================
\section{Introduction}
\label{sec:introduction}

Large Language Models (LLMs) have demonstrated impressive capabilities across diverse domains~\cite{brown2020language,openai2023gpt4}, yet deploying them in safety-critical medical applications reveals fundamental limitations: hallucination of plausible but incorrect facts, lack of grounding in verified clinical knowledge, and inability to reliably perform structured differential diagnosis~\cite{ji2023survey,singhal2023large}.

Retrieval-Augmented Generation (RAG) addresses these limitations by grounding LLM outputs in retrieved evidence~\cite{lewis2020retrieval,gao2023retrieval}. However, standard dense retrieval faces challenges in medical domains: semantic similarity in embedding space does not always capture clinical relevance, and single-source retrieval may miss complementary signals. We propose a \textbf{multi-source RAG framework} that functions as a ``Second Brain'' for LLMs performing medical diagnosis, inspired by Forte's Second Brain methodology~\cite{forte2022building} of capturing and retrieving knowledge through complementary systems. Our architecture integrates:

\begin{enumerate}
    \item \textbf{Dense Vector Retrieval:} Embedding-based similarity search over disease-symptom knowledge, capturing semantic relationships.
    \item \textbf{Sparse Lexical Retrieval with PPMI:} BM25-based retrieval enhanced with PPMI co-occurrence statistics from training data, capturing statistical associations invisible to embeddings.
    \item \textbf{Multi-Source Knowledge Graph:} A graph integrating ontological disease-symptom relationships with data-driven co-occurrence edges weighted by TF-IDF and PPMI scores.
\end{enumerate}

The key insight is that different retrieval modalities capture complementary aspects of medical knowledge. Dense embeddings excel at semantic similarity but may miss statistically significant associations. Sparse retrieval captures these patterns but lacks structural context. Knowledge graphs provide structure but require careful construction. By fusing all three, we achieve near-perfect diagnostic accuracy with zero hallucination.

Our contributions are:
\begin{itemize}
    \item A principled multi-source RAG architecture progressively integrating dense, sparse, and graph-based retrieval (\S\ref{sec:method}).
    \item Empirical demonstration on DDXPlus~\cite{fansi2022ddxplus} and Symptom2Disease~\cite{gretelai_s2d} achieving 98.5\% and 98.8\% Top-5 accuracy (\S\ref{sec:results}).
    \item Evidence that multi-source RAG completely eliminates hallucination (\S\ref{sec:hallucination}).
    \item A model ablation showing RAG compensates for base model differences (\S\ref{sec:model_ablation}).
\end{itemize}

%% ============================================================
\section{Related Work}
\label{sec:related}

\subsection{Retrieval-Augmented Generation}

RAG was introduced by Lewis et al.~\cite{lewis2020retrieval} to augment language model generation with retrieved documents, combining dense passage retrieval~\cite{karpukhin2020dense} with sequence-to-sequence generation. Subsequent work extended RAG to open-domain question answering~\cite{izacard2021leveraging}, dialogue systems~\cite{shuster2021retrieval}, and fact verification~\cite{thorne2018fever}. Recent surveys~\cite{gao2023retrieval,zhao2024retrieval} categorize approaches by retrieval granularity, fusion strategy, and retrieval-generation coupling. Our work contributes to the \emph{multi-source retrieval} dimension, demonstrating that fusing heterogeneous dense, sparse, and graph-based signals provides complementary benefits beyond any single source.

\subsection{Knowledge Graphs in Biomedical AI}

Knowledge graphs structure biomedical knowledge through resources such as UMLS~\cite{bodenreider2004unified}, SNOMED-CT, and domain-specific ontologies~\cite{santos2022knowledge,chandak2023building}. Recent work integrates knowledge graphs with LLMs for medical reasoning~\cite{pan2024unifying}, including graph-enhanced retrieval~\cite{he2024gretriever} and knowledge-grounded generation~\cite{yasunaga2022deep}. Unlike prior approaches that use KGs as static lookup tables or for entity linking, we construct a \emph{multi-source} knowledge graph combining ontological edges with data-driven co-occurrence and PPMI-weighted edges, using it as one component of a broader multi-modal retrieval pipeline.

\subsection{LLM-Based Medical Diagnosis}

Med-PaLM~\cite{singhal2023large} and related systems have shown that LLMs can approach physician-level performance on medical question answering. However, differential diagnosis---ranking multiple possible conditions given a symptom set---remains challenging due to the combinatorial nature of the problem~\cite{mcduff2023towards}. The DDXPlus dataset~\cite{fansi2022ddxplus}, released at NeurIPS 2022, provides a standardized benchmark with 49 pathologies and 223 symptoms.

\subsection{Hallucination Mitigation}

Hallucination---the generation of fluent but factually incorrect content---is particularly dangerous in medical contexts~\cite{ji2023survey,umapathi2023med}. Strategies include retrieval augmentation~\cite{shuster2021retrieval}, constrained decoding~\cite{lu2022neurologic}, and self-consistency checking~\cite{manakul2023selfcheckgpt}. Our work demonstrates that multi-source RAG can \emph{completely eliminate} hallucination in structured diagnostic tasks when the knowledge base covers the target domain.

%% ============================================================
\section{Method}
\label{sec:method}

\subsection{Problem Formulation}

Given a set of patient symptoms $S = \{s_1, s_2, \ldots, s_m\}$ and optional demographic features (age, sex), the task is to produce a ranked list of candidate diagnoses $\hat{D} = [\hat{d}_1, \hat{d}_2, \ldots, \hat{d}_k]$ ordered by predicted likelihood, where the ground-truth diagnosis $d^*$ should appear as high in the ranking as possible. We consider $k=5$ as the primary cutoff.

Performance is measured by:
\begin{itemize}
    \item \textbf{Top-$k$ Accuracy:} $\mathbb{1}[d^* \in \hat{D}_{1:k}]$, averaged over all test cases.
    \item \textbf{NDCG@5:} Normalized Discounted Cumulative Gain at position 5.
    \item \textbf{Macro F1:} Macro-averaged F1 across all disease classes.
    \item \textbf{Hallucination Rate:} $\frac{1}{N}\sum_{i=1}^{N} \mathbb{1}[\exists \hat{d} \in \hat{D}_i : \hat{d} \notin \mathcal{D}]$, where $\mathcal{D}$ is the set of valid diseases.
\end{itemize}

\subsection{Architecture Overview}

Our framework processes each patient case through a pipeline of retrieval modules whose outputs are fused before being presented to the LLM. We evaluate four progressive configurations:

\begin{enumerate}
    \item \textbf{LLM-Only:} Direct prompting with the symptom list.
    \item \textbf{Dense RAG:} LLM augmented with top-$k$ results from dense vector retrieval.
    \item \textbf{BM25+PPMI:} Dense retrieval plus sparse BM25 retrieval with PPMI re-ranking.
    \item \textbf{Multi-Source KG:} Full pipeline adding knowledge graph traversal results.
\end{enumerate}

\begin{figure}[t]
\centering
\fbox{\parbox{0.9\linewidth}{\centering
\textbf{Multi-Source RAG Architecture}\\[6pt]
Patient Symptoms $\rightarrow$ \{Dense Retrieval $\|$ BM25+PPMI $\|$ KG Traversal\}\\
$\downarrow$\\
Reciprocal Rank Fusion\\
$\downarrow$\\
Augmented Prompt $\rightarrow$ LLM $\rightarrow$ Ranked Diagnoses
}}
\caption{Conceptual overview of the multi-source RAG pipeline. Three retrieval modalities provide complementary evidence that is fused via reciprocal rank fusion before augmenting the LLM prompt.}
\label{fig:architecture}
\end{figure}

\subsection{Dense Vector Retrieval}
\label{sec:dense}

Each disease $d_j \in \mathcal{D}$ is represented by a profile document (name, symptoms with prevalence, description) embedded using OpenAI's \texttt{text-embedding-3-small} (1536 dimensions). At inference, the symptom set is embedded as a query and the top-10 most similar profiles retrieved via cosine similarity:
\begin{equation}
    \text{sim}(q, d_j) = \frac{\mathbf{e}_q \cdot \mathbf{e}_{d_j}}{\|\mathbf{e}_q\| \cdot \|\mathbf{e}_{d_j}\|}
\end{equation}
where $\mathbf{e}_q$ and $\mathbf{e}_{d_j} \in \mathbb{R}^{1536}$ are the query and disease profile embeddings.

\subsection{BM25 with PPMI Enhancement}
\label{sec:bm25}

Disease profiles are indexed with BM25~\cite{robertson2009probabilistic}, capturing lexical overlap missed by dense embeddings. We enhance retrieval with PPMI co-occurrence scores from training data. For symptom $s_i$ and disease $d_j$:
\begin{equation}
    \text{PPMI}(s_i, d_j) = \max\left(0, \log_2 \frac{P(s_i, d_j)}{P(s_i) \cdot P(d_j)}\right)
\end{equation}
where $P(s_i, d_j)$ is the joint probability from training data co-occurrence counts and $P(s_i)$, $P(d_j)$ are marginal probabilities. The combined score is:
\begin{equation}
    \text{score}_{\text{BM25+PPMI}}(d_j | S) = \text{BM25}(d_j | S) + \lambda \sum_{s_i \in S} \text{PPMI}(s_i, d_j)
\end{equation}
where $\lambda = 0.5$ balances the BM25 and PPMI components.

\subsection{Multi-Source Knowledge Graph}
\label{sec:kg}

We construct a heterogeneous knowledge graph $G = (V, E)$ with:
\begin{itemize}
    \item \textbf{Nodes} $V = V_D \cup V_S$: Disease nodes $V_D$ and symptom nodes $V_S$.
    \item \textbf{Edges} $E = E_O \cup E_C \cup E_P$: Three types of weighted edges:
    \begin{itemize}
        \item $E_O$ (Ontological): Binary edges from disease definitions indicating known associations.
        \item $E_C$ (Co-occurrence): Edges weighted by normalized co-occurrence frequency $\frac{\text{count}(s_i, d_j)}{\text{count}(d_j)}$ from training data.
        \item $E_P$ (PPMI): Edges weighted by PPMI scores, capturing statistically significant associations beyond base rates.
    \end{itemize}
\end{itemize}

Given query symptoms $S$, we perform a graph traversal starting from symptom nodes $\{v_{s_i} : s_i \in S\}$ and collect neighboring disease nodes. The aggregated relevance score for disease $d_j$ is:
\begin{equation}
    \text{score}_{\text{KG}}(d_j | S) = \sum_{s_i \in S} \left(\alpha \cdot w_O(s_i, d_j) + \beta \cdot w_C(s_i, d_j) + \gamma \cdot w_P(s_i, d_j)\right)
\label{eq:kg_score}
\end{equation}
where $w_O$, $w_C$, $w_P$ are the edge weights for each type (zero if no edge exists), and $\alpha = \beta = \gamma = 1/3$ are equal fusion coefficients.

\subsection{Retrieval Fusion and LLM Prompting}
\label{sec:fusion}

Retrieved candidates from all active retrieval layers are merged using Reciprocal Rank Fusion (RRF)~\cite{cormack2009reciprocal}:
\begin{equation}
    \text{RRF}(d_j) = \sum_{r \in R} \frac{1}{k_{\text{rrf}} + \text{rank}_r(d_j)}
\end{equation}
where $R$ is the set of active retrieval sources, $\text{rank}_r(d_j)$ is the rank of disease $d_j$ in source $r$'s results, and $k_{\text{rrf}} = 60$ is a smoothing constant.

The top-ranked disease profiles (up to 10) are formatted into structured context and included in the LLM prompt. The prompt instructs the LLM to consider only the provided candidate diseases, rank them by likelihood given the patient's symptoms, and output a structured JSON list. This constrained generation approach ensures that the LLM's output is grounded in retrieved evidence, directly mitigating hallucination.

%% ============================================================
\section{Experimental Setup}
\label{sec:experiments}

\subsection{Datasets}

\paragraph{DDXPlus~\cite{fansi2022ddxplus}.} A large-scale dataset containing ${\sim}$1.3M patient cases across 49 pathologies and 223 symptoms, constructed from a validated clinical knowledge base. We evaluate on $n{=}200$ stratified samples from the 134K test set, ensuring proportional disease representation. Each case includes binary/categorical symptom features, demographics, and a ground-truth differential diagnosis. The KG and PPMI statistics are built from 10K training cases.

\paragraph{Symptom2Disease~\cite{gretelai_s2d}.} A dataset of ${\sim}$1,065 free-text symptom descriptions mapped to 24 diseases. Unlike DDXPlus, symptoms are natural language rather than structured features. We split 70/30 into training ($n{=}745$) and test ($n{=}320$), constructing a \emph{domain-specific} KG entirely from S2D training data to test cross-domain generalization.

\subsection{Model Configuration}

GPT-4o-mini (\texttt{gpt-4o-mini-2024-07-18}) serves as primary LLM with temperature 0.0 and seed 42. GPT-4o (\texttt{gpt-4o-2024-08-06}) is used for model ablation. Dense embeddings use \texttt{text-embedding-3-small}. All experiments use the same prompt template, differing only in retrieved context. Hallucination is measured by checking whether predicted diseases appear in the target dataset's vocabulary.

%% ============================================================
\section{Results}
\label{sec:results}

\subsection{DDXPlus Benchmark Results}

Table~\ref{tab:ddxplus} presents results on the DDXPlus dataset using GPT-4o-mini. The unaugmented LLM achieves only 32.5\% Top-1 accuracy with a 59.8\% hallucination rate, underscoring the inadequacy of relying solely on parametric knowledge for medical diagnosis.

\begin{table}[t]
\centering
\caption{Results on DDXPlus ($n{=}200$, 49 diseases, 223 symptoms) using GPT-4o-mini. Best results in \textbf{bold}.}
\label{tab:ddxplus}
\begin{tabular}{lcccccc}
\toprule
\textbf{Method} & \textbf{Top-1} & \textbf{Top-3} & \textbf{Top-5} & \textbf{NDCG@5} & \textbf{F1} & \textbf{Halluc.} \\
\midrule
LLM-Only       & 0.325 & 0.420 & 0.485 & 0.415 & 0.163 & 0.598 \\
Dense RAG       & 0.855 & 0.935 & 0.955 & 0.934 & 0.321 & \textbf{0.000} \\
BM25+PPMI       & 0.955 & 0.965 & 0.965 & 0.993 & 0.325 & 0.008 \\
Multi-Source KG & \textbf{0.955} & \textbf{0.985} & \textbf{0.985} & \textbf{0.999} & \textbf{0.331} & \textbf{0.000} \\
\bottomrule
\end{tabular}
\end{table}

Dense RAG improves to 85.5\% Top-1 and eliminates hallucination. BM25+PPMI further improves Top-1 to 95.5\%---a 10pp gain---capturing complementary signals missed by embeddings. The full Multi-Source KG matches 95.5\% Top-1 but improves Top-3 to 98.5\% and NDCG@5 to 0.999, indicating that the KG primarily improves \emph{ranking quality}. It also eliminates the residual 0.8\% hallucination from BM25+PPMI.

\subsection{Model Ablation: GPT-4o vs.\ GPT-4o-mini}
\label{sec:model_ablation}

To assess whether our framework's effectiveness depends on the specific base model, we replicate the DDXPlus experiment using GPT-4o. Table~\ref{tab:gpt4o} presents the results.

\begin{table}[t]
\centering
\caption{Model ablation on DDXPlus ($n{=}200$) using GPT-4o. Compare with Table~\ref{tab:ddxplus} (GPT-4o-mini).}
\label{tab:gpt4o}
\begin{tabular}{lcccccc}
\toprule
\textbf{Method} & \textbf{Top-1} & \textbf{Top-3} & \textbf{Top-5} & \textbf{NDCG@5} & \textbf{F1} & \textbf{Halluc.} \\
\midrule
LLM-Only       & 0.235 & 0.365 & 0.430 & 0.356 & 0.151 & 0.672 \\
Dense RAG       & 0.820 & 0.925 & 0.955 & 0.918 & 0.321 & 0.002 \\
BM25+PPMI       & 0.960 & 0.970 & 0.970 & 0.992 & 0.326 & 0.013 \\
Multi-Source KG & 0.955 & 0.985 & 0.985 & 0.999 & 0.331 & 0.003 \\
\bottomrule
\end{tabular}
\end{table}

GPT-4o performs \emph{worse} without RAG (23.5\% vs.\ 32.5\% Top-1, 67.2\% vs.\ 59.8\% hallucination), suggesting the larger model's broader knowledge introduces more spurious associations when unconstrained. However, with full RAG, both models converge to identical performance (95.5\% Top-1, 98.5\% Top-5, NDCG@5 of 0.999), validating that \textbf{the retrieval framework compensates for variation in base model capability}. This implies that retrieval infrastructure yields greater returns than larger base models for structured diagnostic tasks.

\subsection{Cross-Domain Results: Symptom2Disease}
\label{sec:cross_domain}

To evaluate generalization beyond DDXPlus, we construct a domain-specific knowledge graph from the Symptom2Disease training split and evaluate on the held-out test set. Table~\ref{tab:s2d} presents the results.

\begin{table}[t]
\centering
\caption{Results on Symptom2Disease ($n{=}320$, 24 diseases) using GPT-4o-mini with a domain-specific knowledge graph built from the S2D training split.}
\label{tab:s2d}
\begin{tabular}{lcccccc}
\toprule
\textbf{Method} & \textbf{Top-1} & \textbf{Top-3} & \textbf{Top-5} & \textbf{NDCG@5} & \textbf{F1} & \textbf{Halluc.} \\
\midrule
LLM-Only       & 0.325 & 0.478 & 0.569 & 0.494 & 0.196 & 0.759 \\
Dense RAG       & 0.800 & 0.916 & 0.953 & 0.883 & 0.318 & \textbf{0.000} \\
BM25+PPMI       & 0.766 & 0.794 & 0.794 & 0.784 & 0.265 & 0.004 \\
Multi-Source KG & \textbf{0.903} & \textbf{0.978} & \textbf{0.988} & \textbf{0.952} & \textbf{0.329} & \textbf{0.000} \\
\bottomrule
\end{tabular}
\end{table}

The Multi-Source KG achieves 90.3\% Top-1 and 98.8\% Top-5 with zero hallucinations, versus 32.5\% Top-1 and 75.9\% hallucination for LLM-Only, confirming cross-domain generalization.

Notably, BM25+PPMI \emph{underperforms} Dense RAG on S2D (76.6\% vs.\ 80.0\% Top-1)---a reversal from DDXPlus---because BM25's lexical matching is less effective for S2D's varied natural language expressions versus structured symptom codes. The Multi-Source KG provides its largest relative improvement here, boosting Top-1 from 80.0\% to 90.3\% (+10.3pp), demonstrating that the KG is especially valuable when individual modalities show mixed performance.

\subsection{Ablation Analysis}
\label{sec:ablation}

Table~\ref{tab:ablation} quantifies the incremental contribution of each retrieval layer on DDXPlus.

\begin{table}[t]
\centering
\caption{Ablation: incremental contribution of each retrieval layer on DDXPlus.}
\label{tab:ablation}
\begin{tabular}{lccc}
\toprule
\textbf{Transition} & \textbf{$\Delta$Top-1} & \textbf{$\Delta$Top-5} & \textbf{$\Delta$NDCG@5} \\
\midrule
Baseline $\to$ +Dense RAG    & +0.530 & +0.470 & +0.519 \\
+Dense $\to$ +BM25+PPMI      & +0.100 & +0.010 & +0.059 \\
+BM25+PPMI $\to$ +KG         & +0.000 & +0.020 & +0.006 \\
\midrule
\textbf{Total improvement}   & \textbf{+0.630} & \textbf{+0.500} & \textbf{+0.584} \\
\bottomrule
\end{tabular}
\end{table}

\textbf{Dense RAG} provides the largest single improvement (+53.0 pp Top-1), transforming the LLM from an unreliable guesser (32.5\% accuracy, 59.8\% hallucination) into a competent diagnostic assistant (85.5\% accuracy, 0\% hallucination). \textbf{BM25+PPMI} adds a substantial 10.0 pp Top-1 improvement by capturing statistical co-occurrence patterns that dense embeddings miss---symptoms that rarely appear in the corpus but are highly specific to certain diseases may not cluster closely in embedding space yet provide strong diagnostic signal. \textbf{Multi-Source KG} provides improvement primarily in ranking quality: NDCG@5 reaches 0.999 (near-perfect), and Top-3/Top-5 improve by 2.0 pp each. While the marginal gain is modest on DDXPlus, the KG's contribution is more pronounced on S2D (Section~\ref{sec:cross_domain}), where it provides the largest relative improvement of any single component.

\subsection{Hallucination Analysis}
\label{sec:hallucination}

The hallucination results deserve special attention given the safety-critical nature of medical diagnosis (Table~\ref{tab:halluc}).

\begin{table}[t]
\centering
\caption{Hallucination rates across methods, datasets, and models.}
\label{tab:halluc}
\begin{tabular}{lccc}
\toprule
\textbf{Method} & \textbf{DDXPlus (4o-mini)} & \textbf{DDXPlus (4o)} & \textbf{S2D (4o-mini)} \\
\midrule
LLM-Only        & 0.598 & 0.672 & 0.759 \\
Dense RAG        & 0.000 & 0.002 & 0.000 \\
BM25+PPMI        & 0.008 & 0.013 & 0.004 \\
Multi-Source KG  & 0.000 & 0.003 & 0.000 \\
\bottomrule
\end{tabular}
\end{table}

Without RAG, the LLM hallucinates in 60--76\% of cases. This rate is \emph{higher} for GPT-4o (67.2\%) than GPT-4o-mini (59.8\%), suggesting that larger models may generate more diverse but invalid disease names. S2D shows the highest baseline hallucination (75.9\%), reflecting the model's vocabulary exceeding the dataset's 24-disease scope. With the full Multi-Source KG pipeline, hallucination is eliminated entirely on both datasets with GPT-4o-mini (0.000) and reduced to a negligible 0.3\% with GPT-4o.

%% ============================================================
\section{Discussion}
\label{sec:discussion}

\subsection{The Progressive Value of Multi-Source Retrieval}

Our results demonstrate clear progressive improvement from each retrieval modality, supporting the multi-source ``Second Brain'' design principle. Dense embeddings capture broad semantic similarity, BM25+PPMI captures statistical co-occurrence patterns, and the knowledge graph captures structural ontological relationships weighted by multiple evidence sources. This complementarity mirrors how human clinicians reason: combining textbook knowledge (ontological), clinical experience (co-occurrence), and pattern matching (semantic similarity).

The S2D results further illuminate this interplay. When sparse retrieval underperforms dense retrieval (due to free-text vs.\ structured inputs), the Multi-Source KG compensates by integrating both signals, achieving best performance regardless of input format.

\subsection{Model-Agnostic Retrieval: RAG as an Equalizer}

The GPT-4o ablation reveals that RAG augmentation acts as an \emph{equalizer} across model capabilities. The 9.0-percentage-point gap in the LLM-Only condition vanishes entirely with Multi-Source KG augmentation. This suggests that for well-defined diagnostic tasks: (1) the base model's parametric medical knowledge matters less than the quality of retrieved evidence; (2) smaller, cheaper models can match larger models with appropriate retrieval infrastructure; and (3) the retrieval framework provides a performance floor largely independent of model size. Organizations can deploy cost-effective models without sacrificing diagnostic accuracy.

\subsection{Domain Generalization and Practical Implications}

The strong S2D results demonstrate generalization when domain-specific KGs are constructed automatically from training data, requiring no expert curation. The performance gap between DDXPlus (95.5\%) and S2D (90.3\%) reflects genuine task difficulty: S2D's free-text descriptions have higher linguistic variability and fewer training examples per disease (${\sim}$31 vs.\ ${\sim}$204). Hybrid approaches combining automatic construction with expert curation could further improve performance.

Our use of TF-IDF and PPMI as clinical significance proxies is effective but may not capture pathognomonic symptoms, severity weighting, or temporal patterns. The near-perfect performance (98.5--98.8\% Top-5, zero hallucination) suggests viability as clinical decision support, with considerations for knowledge base maintenance, scope transparency, domain adaptation, and regulatory compliance.

%% ============================================================
\section{Limitations}
\label{sec:limitations}

We acknowledge several limitations. DDXPlus covers 49 diseases and S2D covers 24---both subsets of real clinical pathology. DDXPlus cases are synthetically generated and may not capture the full complexity of real presentations, including comorbidities and atypical cases. Our evaluation uses $n{=}200$ and $n{=}320$ test cases; while sufficient for large effect sizes, larger-scale validation is warranted. We evaluate only GPT-4o-mini and GPT-4o; interactions with open-source medical models (e.g., Meditron~\cite{chen2023meditron}) may differ. Our static KGs do not incorporate external medical knowledge bases (UMLS, SNOMED-CT) or adapt over time, and we use equal fusion weights ($\alpha{=}\beta{=}\gamma{=}1/3$) rather than learned weights---though near-perfect results leave limited room for improvement.

%% ============================================================
\section{Conclusion}
\label{sec:conclusion}

We presented a multi-source RAG framework serving as a ``Second Brain'' for LLM-based medical diagnosis. By integrating dense vector retrieval, BM25 with PPMI statistics, and a multi-source knowledge graph, our system achieves 95.5\% Top-1 and 98.5\% Top-5 on DDXPlus, and 90.3\% Top-1 and 98.8\% Top-5 on Symptom2Disease---with zero hallucinations---versus baselines of 32.5\% Top-1 and 60--76\% hallucination. Each layer contributes complementary information, a GPT-4o ablation confirms model-agnostic effectiveness, and the pipeline generalizes via automatic KG construction. This paradigm offers a principled path toward reliable clinical decision support---not as an autonomous diagnostician, but as a grounded ``Second Brain'' augmenting clinical reasoning with structured, multi-modal medical knowledge.

%% ============================================================
\section*{CRediT Authorship Contribution Statement}

\textbf{First Author:} Conceptualization, Methodology, Software, Validation, Writing -- Original Draft.
\textbf{Second Author:} Methodology, Formal Analysis, Writing -- Review \& Editing.
\textbf{Third Author:} Supervision, Writing -- Review \& Editing.

\section*{Declaration of Competing Interest}

The authors declare that they have no known competing financial interests or personal relationships that could have appeared to influence the work reported in this paper.

\section*{Data Availability}

The DDXPlus dataset is publicly available at \url{https://github.com/mila-iqia/ddxplus}. The Symptom2Disease dataset is available at \url{https://huggingface.co/datasets/gretelai/symptom_to_diagnosis}. Code and knowledge graph construction scripts will be released upon publication.

%% ============================================================
\bibliographystyle{elsarticle-num}

\begin{thebibliography}{27}

\bibitem{brown2020language}
T.~Brown, B.~Mann, N.~Ryder, et~al., Language models are few-shot learners, in: Advances in Neural Information Processing Systems, vol.~33, 2020, pp.~1877--1901.

\bibitem{openai2023gpt4}
OpenAI, GPT-4 technical report, arXiv preprint arXiv:2303.08774, 2023.

\bibitem{ji2023survey}
Z.~Ji, N.~Lee, R.~Frieske, et~al., Survey of hallucination in natural language generation, ACM Computing Surveys 55~(12) (2023) 1--38.

\bibitem{singhal2023large}
K.~Singhal, S.~Azizi, T.~Tu, et~al., Large language models encode clinical knowledge, Nature 620~(7972) (2023) 172--180.

\bibitem{lewis2020retrieval}
P.~Lewis, E.~Perez, A.~Piktus, et~al., Retrieval-augmented generation for knowledge-intensive NLP tasks, in: Advances in Neural Information Processing Systems, vol.~33, 2020, pp.~9459--9474.

\bibitem{gao2023retrieval}
Y.~Gao, Y.~Xiong, A.~Gupta, et~al., Retrieval-augmented generation for large language models: A survey, arXiv preprint arXiv:2312.10997, 2023.

\bibitem{forte2022building}
T.~Forte, Building a Second Brain: A Proven Method to Organize Your Digital Life and Unlock Your Creative Potential, Atria Books, 2022.

\bibitem{robertson2009probabilistic}
S.~Robertson, H.~Zaragoza, The probabilistic relevance framework: BM25 and beyond, Foundations and Trends in Information Retrieval 3~(4) (2009) 333--389.

\bibitem{fansi2022ddxplus}
A.~Fansi~Tchango, R.~Camara, I.~Kreitchmann, et~al., DDXPlus: A new dataset for automatic medical diagnosis, in: Advances in Neural Information Processing Systems, vol.~35, 2022.

\bibitem{gretelai_s2d}
Gretel.ai, Symptom to diagnosis dataset, Hugging Face Datasets, 2023. URL: \url{https://huggingface.co/datasets/gretelai/symptom_to_diagnosis}.

\bibitem{santos2022knowledge}
A.~Santos, A.~Coelho, et~al., A knowledge graph to interpret clinical proteomics data, Nature Biotechnology 40 (2022) 692--702.

\bibitem{chandak2023building}
P.~Chandak, K.~Huang, M.~Zitnik, Building a knowledge graph to enable precision medicine, Scientific Data 10 (2023) 67.

\bibitem{bodenreider2004unified}
O.~Bodenreider, The Unified Medical Language System (UMLS): Integrating biomedical terminology, Nucleic Acids Research 32 (2004) D267--D270.

\bibitem{pan2024unifying}
S.~Pan, L.~Luo, Y.~Wang, et~al., Unifying large language models and knowledge graphs: A roadmap, IEEE Transactions on Knowledge and Data Engineering (2024).

\bibitem{mcduff2023towards}
D.~McDuff, M.~Schaekermann, T.~Tu, et~al., Towards accurate differential diagnosis with large language models, arXiv preprint arXiv:2312.00164, 2023.

\bibitem{karpukhin2020dense}
V.~Karpukhin, B.~O\u{g}uz, S.~Min, et~al., Dense passage retrieval for open-domain question answering, in: EMNLP, 2020, pp.~6769--6781.

\bibitem{izacard2021leveraging}
G.~Izacard, E.~Grave, Leveraging passage retrieval with generative models for open domain question answering, in: EACL, 2021, pp.~874--880.

\bibitem{shuster2021retrieval}
K.~Shuster, S.~Poff, M.~Chen, et~al., Retrieval augmentation reduces hallucination in conversation, in: EMNLP Findings, 2021.

\bibitem{thorne2018fever}
J.~Thorne, A.~Vlachos, C.~Christodoulopoulos, A.~Mittal, FEVER: A large-scale dataset for fact extraction and verification, in: NAACL-HLT, 2018, pp.~809--819.

\bibitem{zhao2024retrieval}
P.~Zhao, H.~Zhang, Q.~Yu, et~al., Retrieval-augmented generation for AI-generated content: A survey, arXiv preprint arXiv:2402.19473, 2024.

\bibitem{umapathi2023med}
L.~K.~Umapathi, A.~Pal, M.~Sankarasubbu, Med-HALT: Medical domain hallucination test for large language models, in: CoNLL, 2023.

\bibitem{lu2022neurologic}
X.~Lu, S.~Welleck, P.~West, et~al., NeuroLogic A*esque decoding: Constrained text generation with lookahead heuristics, in: NAACL, 2022.

\bibitem{manakul2023selfcheckgpt}
P.~Manakul, A.~Liusie, M.~J.~F.~Gales, SelfCheckGPT: Zero-resource black-box hallucination detection for generative large language models, in: EMNLP, 2023.

\bibitem{cormack2009reciprocal}
G.~V.~Cormack, C.~L.~A.~Clarke, S.~B\"uttcher, Reciprocal rank fusion outperforms condorcet and individual rank learning methods, in: SIGIR, 2009, pp.~758--759.

\bibitem{he2024gretriever}
X.~He, Y.~Tian, Y.~Sun, et~al., G-Retriever: Retrieval-augmented generation for textual graph understanding and question answering, arXiv preprint arXiv:2402.07630, 2024.

\bibitem{yasunaga2022deep}
M.~Yasunaga, A.~Bosselut, H.~Ren, et~al., Deep bidirectional language-knowledge graph pretraining, in: NeurIPS, 2022.

\bibitem{chen2023meditron}
Z.~Chen, A.~Cano, et~al., Meditron-70B: Scaling medical pretraining for large language models, arXiv preprint arXiv:2311.16079, 2023.

\end{thebibliography}

\end{document}
